\chapter{Justification et argumentaire}
\label{chap:justification}

\begin{todobox}[Chapitre à compléter après chiffrage]
Ce chapitre est destiné à la direction. Il présente les scénarios d'investissement, les coûts et les bénéfices.
\end{todobox}

\section{Scénarios d'investissement}
\begin{todobox}[À produire après audit]
\begin{itemize}
  \item Scénario minimal : RAM + disques R630 uniquement.
  \item Scénario cible : RAM + disques + réseau 10 GbE + éventuellement serveurs complémentaires.
\end{itemize}
\end{todobox}

\section{Argumentaire pour la direction}
\begin{itemize}
  \item La plateforme historique a déjà supporté 500 utilisateurs. Sa remise en service est prioritaire.
  \item La phase transitoire permet de fournir des services rapidement sans attendre la reconstruction.
  \item L'architecture est modulaire : chaque investissement débloque une étape.
  \item L'infrastructure est exploitable par une personne formée, sans dépendance à un prestataire.
  \item La trajectoire OpenStack aligne l'école sur les standards de l'industrie du cloud.
\end{itemize}

\section{Risques et mitigation}
\begin{tabularx}{\textwidth}{@{}lXX@{}}
\toprule
\textbf{Risque} & \textbf{Impact} & \textbf{Mitigation} \\
\midrule
Matériel historique non réutilisable & Reconstruction complète = coût élevé & Audit préalable, scénarios alternatifs \\
Perte de compétence (départ admin) & Infrastructure non maintenue & Documentation, formation, externalisation minimale \\
Dépassement budgétaire & Retard de phases & Découpage en lots, jalons de validation \\
Attaque cyberlab & Compromission de l'infra & Isolation stricte, supervision \\
SRX300 sous-dimensionné & VPN étudiant impossible & Concentrateur VPN dédié, portail web \\
\bottomrule
\end{tabularx}